% =====================================================================
%  Aufgabenstellung (konsolidiert) — DEUTSCH (Original / Hauptsprache)
%  Allgemeiner Scope (Suchalgorithmen/Suchbäume), design-space-Vokabular
%  mit Zitat; KEINE Achsen-Liste, KEINE Tier/Organ-Metapher (das ist Lösung,
%  linearer Aufbau). Quellen: docs/termine (Termin 1-9) + Idreos design space.
%  Englisches Pendant: en.tex.
% =====================================================================
{\large\textbf{Aufgabenstellung}}\\[0.4em]
{\itshape Aktive cache-bewusste Hardware-Adaption:\\ Eine Cache-Engine für Trie-basierte Indexstrukturen}

\bigskip
\minisec{Kontext und Motivation}
Suchalgorithmen und Indexstrukturen sind ein Kernbaustein moderner
In-Memory-Datenbanksysteme. Die Professur für Datenbanken an der TU Dresden adressiert
die hardware-nahe Optimierung solcher Systeme. Seit der Einführung des Adaptive
Radix Tree wurde eine Vielzahl cache-bewusster Techniken und verwandter
Algorithmenfamilien vorgeschlagen --- von Knoten-Layout über Prefetching und
Succinct-Repräsentationen bis zur Allokation ---, doch nutzt jede Veröffentlichung eigene Datenstrukturen, Lastprofile und
Messmethodik. Ergebnisse sind dadurch über Veröffentlichungen hinweg schwer
vergleichbar.

\minisec{Problemstellung}
Suchverfahren --- von ausgewogenen Suchbäumen (B-/B\textsuperscript{+}-Bäume)
über Tries und Radix-Bäume bis zu spezialisierten, eigenentwickelten Konstrukten
--- spannen einen großen \emph{Entwurfsraum} (\emph{design space}) auf:
Bereits eine zwei- bzw.\ dreikomponentige Datenstruktur erlaubt astronomisch
viele Entwürfe, und ein konkreter Algorithmus ist ein einzelner Punkt in diesem
Raum, beschrieben durch viele einzelne Entwurfsentscheidungen --- wie Knoten
repräsentiert, durchlaufen, im Speicher angeordnet und alloziert werden
\cite{idreos2018periodic, idreos2018datacalculator}. Cache-line-bewusstes
Verhalten entsteht dabei nicht an einer einzelnen dieser Stellen, sondern an
\emph{vielen verschiedenen strukturellen Stellen} eines Algorithmus und aus deren
Zusammenspiel. Bestehende Arbeiten optimieren meist einzelne Stellen isoliert und
berichten schwer vergleichbare Ergebnisse. Es fehlt ein gemeinsames
\textbf{Prinzip}, das cache-line-bewusstes Verhalten an den verschiedenen Stellen
eines Suchalgorithmus --- austauschbar, systematisch und fair --- messen und
auswerten kann.

\minisec{Zielsetzung}
Ziel der Arbeit ist es, ein solches Prinzip (ein Rahmenwerk) zu finden, zu
entwerfen und prototypisch umzusetzen, das (i) cache-bewusste Suchalgorithmen in
austauschbare Entwurfsbestandteile zerlegt, (ii) Algorithmen des Stands der
Technik als explizite Konfigurationen (Punkte im Entwurfsraum) rekonstruiert und
(iii) cache-line-bewusstes Verhalten sowohl an den einzelnen Stellen als auch für
den gesamten Algorithmus auf gemeinsamer CPU-only-Basis misst und auswertet. Als
Prüfling für die Zugehörigkeit zum Stand der Technik trägt die Arbeit
\mbox{PRT-ART} bei (einen Permutations-Trie-/ART-Hybriden).

\minisec{Teilaufgaben}
\begin{enumerate}
  \item Systematische Erfassung der Entwurfsbestandteile cache-bewusster
        Suchalgorithmen und Identifikation der Stellen, an denen
        cache-line-bewusste Entscheidungen wirksam werden.
  \item Rekonstruktion von Entwürfen des Stands der Technik (ART, HOT, Masstree,
        CoCo-trie, START, B\textsuperscript{2}-Tree, Wormhole, SuRF sowie weitere
        Verfahren) als explizite, vergleichbare Konfigurationen.
  \item Entwurf und Implementierung von PRT-ART als Prüfling.
  \item Aufbau einer reproduzierbaren Mess-Pipeline
        (Binary~$\to$~CSV~$\to$~LaTeX~$\to$~Diagramm) mit profilbewusstem
        YCSB-Workload-Routing und Hardware-Counter-Erfassung.
  \item Systematische, möglichst vollständige Messung über die Konfigurationen in
        drei verpflichtenden Messreihen: (A)~Prüfling gegen den Stand der Technik,
        (B)~systematische Variation der Entwurfsentscheidungen,
        (C)~Merge/Regression alter gegen neue Varianten.
  \item Auswertung: an welchen Stellen und in welchen Kombinationen
        cache-line-bewusste Entscheidungen für welche datentyp-spezifischen
        Lastmuster vorteilhaft oder nachteilig sind.
\end{enumerate}

\minisec{Methodik und Evaluation}
Die Evaluation ist CPU-only und vergleicht die Konfigurationen gegen etablierte
Baselines unter variierten Datenvolumina, Keylängen, Präfixüberlappungen, lokaler
Dichte, Skew, Treffer-/Fehlerquoten und Value-Größen. Gemessen werden
Exact-/Prefix-Lookup, Prefix-Enumeration, Insert/Update/Delete,
p50/p95/p99-Latenz, Durchsatz, Bytes pro Key sowie Hardwarezähler (cycles,
instructions, L1/L2/LLC- und dTLB-Misses, Branch-Mispredictions). Zunächst wird
Single-Thread-Verhalten untersucht, danach folgen leseparallele
Multi-Thread-Szenarien.

\minisec{Umfangsabgrenzung}
Der Kernbeitrag ist CPU-only und single-threaded mit Lese-Skalierung. GPU liegt
außerhalb des Umfangs; Schreibparallelität, SIMD und Engine-/Optimierer-Integration
sind optionale Erweiterungen.
